\subsection{Capture-Availability Pacing}
\label{sec:pacing}

\emph{Pacing} maintains a running forecast of draft workload to regulate the next capture opportunity.
It extends this forecast after each decision and rebases it as draft sequences complete.

\smallskip\noindent\textbf{Delay rule.}
When the \texttt{captureAvailable} callback for capture \(i\) is ready for delivery, \(\mathit{controller}.\textsc{DecidePacingDelay}\) (Algorithm~\ref{alg:overview}, line~\ref{alg:ov:pacing}) determines the callback delivery delay.
Pacing returns zero delay if the draft queue is empty or no draft sequence has completed.
Otherwise, it applies
\begin{equation*}
  \mathit{delay}_i
  =
  \left\lceil
    \frac{
      \max\{0,\hat B_i+2\hat R-T_{\mathrm{queue}}\}
    }{2}
  \right\rceil.
\end{equation*}

First, the predicted backlog \(\hat B_i\) estimates the time needed to drain draft work from earlier captures.
At each pacing decision, pacing computes this backlog as
\begin{equation*}
  \hat B_i
  =
  \max\{0,\hat t_{\mathrm{drain}}-\mathit{now}\},
\end{equation*}
where \(\hat t_{\mathrm{drain}}\) is the predicted drain time, computed from point estimates of draft work and updated by Equations~\eqref{eq:clock-update} and~\eqref{eq:clock-rebase}.

Second, pacing uses a conservative duration reserve \(\hat R\) per capture to project the draft workload added by the current capture \(i\) and the next capture \(i+1\) enabled by this callback.
Because \(\hat B_i\) does not yet include draft work for either capture, pacing adds \(\hat R\) for each, giving the projected workload \(\hat B_i+2\hat R\).
Pacing computes the reserve as
\begin{equation*}
  \hat R
  =
  \hat P(\mathcal K^{\mathrm{adm}})
  +
  \max\!\left(0,C_\tau-\hat P(\mathcal K)\right),
\end{equation*}
where \(\mathcal K^{\mathrm{adm}}\) is the subsequence of the configured sequence \(\mathcal K\) excluding optional work disabled by admission, and \(C_\tau\) is the recency-weighted \(\tau\)-quantile of completed draft-sequence durations.
Pacing applies the same recency weighting and selects \(\tau\) using the same effective-sample-size-based rule as admission does for residual factors in Section~\ref{sec:admission}.
The second term \(\max(0,C_\tau-\hat P(\mathcal K))\) adds empirical headroom because the draft sequences of captures \(i\) and \(i+1\) execute after the existing backlog and may encounter further thermal throttling.
With all configured stages enabled (\(\mathcal K^{\mathrm{adm}}=\mathcal K\)), pacing uses the larger of \(C_\tau\) and \(\hat P(\mathcal K)\) as its reserve.
After admission disables optional work, pacing adds the positive excess of \(C_\tau\) over \(\hat P(\mathcal K)\) as headroom to the point estimate \(\hat P(\mathcal K^{\mathrm{adm}})\).

Third, \(T_{\mathrm{queue}}\) denotes the remaining timeout budget of the most recently enqueued capture in the draft queue.
At the pacing decision point, capture \(i+1\)'s timeout clock has not yet started.
\(T_{\mathrm{queue}}\) therefore serves as a queue-level pressure reference for the projected workload \(\hat B_i+2\hat R\).
Pacing uses this pressure to limit queue growth, aiming to preserve time for capture \(i+1\)'s draft processing after acceptance.

Finally, pacing sets the callback delivery delay to half of the projected workload excess over the queue-level budget.
This heuristic aims to limit irreversible shot-to-shot latency, leaving workload reduction to admission.

\smallskip\noindent\textbf{Backlog extension.}
In the parallel capture workflow, the pacing decision for capture \(i+1\) may occur before capture \(i\)'s draft sequence is enqueued.
Pacing therefore provisionally extends \(\hat t_{\mathrm{drain}}\) after computing \(\mathit{delay}_i\) to include this work in subsequent delay calculations.
It uses \(\hat P^{\mathrm{last}}\), the value of \(\hat P(\mathcal K^{\mathrm{adm}})\) evaluated at the most recent draft completion:
\begin{equation}
  \hat t_{\mathrm{drain}}
  \gets
  \max\{\hat t_{\mathrm{drain}},\mathit{now}+\mathit{delay}_i\}
  +\lceil\hat P^{\mathrm{last}}\rceil.
  \label{eq:clock-update}
\end{equation}
The \(\max\) term anchors this provisional forecast to the later of the existing drain forecast \(\hat t_{\mathrm{drain}}\) and the callback delivery time \(\mathit{now}+\mathit{delay}_i\).
This keeps at least one draft sequence's predicted processing time in the provisional backlog when the next capture opportunity opens.
In this way, this extension accounts for the possibility that capture \(i\)'s draft sequence has not started by the callback delivery time.

\smallskip\noindent\textbf{Backlog rebase.}
Pacing rebases the backlog at each draft completion because provisional backlog extensions use earlier point estimates that can become stale as durations and admission state change.
After the shared predictor and admission update their states (Algorithm~\ref{alg:overview}, line~\ref{alg:ov:update}), pacing re-anchors the drain forecast at \(\mathit{now}\):
\begin{equation}
  \hat t_{\mathrm{drain}}
  \gets
  \mathit{now}
  +\left\lceil
    N_{\mathrm{pend}}\cdot\hat P(\mathcal K^{\mathrm{adm}})
  \right\rceil.
  \label{eq:clock-rebase}
\end{equation}
Here, \(N_{\mathrm{pend}}\) is the number of captures whose pacing decisions have occurred but whose draft sequences have not yet started.
Pacing increments this count after each pacing decision and decrements it when the corresponding draft sequence starts.
All pending captures share the configured sequence \(\mathcal K\), which cannot change until the draft queue drains.
Backlog rebase therefore refreshes the forecast for all pending captures using the latest duration estimates and admission state.
